% ==============================================================================
% Plantilla Institucional de Trabajo Terminal — UPIIZ IPN
% Archivo: capitulos/05-estado-del-arte.tex
% ==============================================================================

\chapter{Estado del arte}
\label{cap:estado-del-arte}

% ==============================================================================
% REQUISITO INSTITUCIONAL:
% Extensión: Máximo CUATRO páginas.
%
% INSTRUCCIONES PEDAGÓGICAS:
% Realice una revisión crítica y comparativa de las tecnologías, productos comerciales,
% patentes y publicaciones científicas recientes (últimos 5 años preferentemente)
% relacionadas directamente con el problema que resuelve su proyecto.
%
% No se limite a listar artículos; elabore un análisis comparativo destacando:
% - Ventajas y desventajas de cada enfoque existente.
% - Parámetros de desempeño, costos, complejidad y limitaciones.
% - La brecha tecnológica u oportunidad que justifica la solución propuesta en este TT.
% ==============================================================================

\section{Tecnologías y soluciones existentes}
\label{sec:arte-tecnologias}

% Escriba aquí la revisión de soluciones comerciales y académicas...
En el ámbito internacional, el desarrollo de soluciones mecatrónicas para este tipo de aplicaciones ha experimentado un crecimiento notable impulsado por avances en microelectrónica y algoritmos de optimización. Diversos autores han propuesto arquitecturas basadas en controladores lógicos programables, microcontroladores de 32 bits y computadoras de placa reducida (SBC) \cite{ogata2010modern, arizona2023embedded}.

Asimismo, en la literatura especializada se reportan implementaciones con diferentes grados de autonomía y precisión, empleando esquemas de comunicación inalámbrica y protocolos industriales \cite{smith2022control, ieee2021robotics}.

\section{Comparativa técnica de enfoques}
\label{sec:arte-comparativa}

A fin de situar la presente propuesta frente a las alternativas documentadas, la \tabref{tab:comparativa-arte} resume las principales características operativas de las tecnologías analizadas.

\begin{table}[htbp]
    \centering
    \caption{Comparación técnica de tecnologías y enfoques existentes.}
    \label{tab:comparativa-arte}
    \begin{tabularx}{\textwidth}{p{3.0cm} X X p{3.2cm}}
        \toprule
        \textbf{Enfoque / Solución} & \textbf{Ventajas} & \textbf{Desventajas} & \textbf{Aplicabilidad en este TT} \\
        \midrule
        Sistemas comerciales propietarios & Alta robustez y soporte de fábrica & Costo elevado y nula personalización & Limitada para fines académicos \\
        Plataformas abiertas modulares & Bajo costo y alta flexibilidad & Mayor complejidad de integración & Base fundamental de este trabajo \\
        Control clásico analógico & Simplicidad circuital & Sensibilidad a ruido y deriva térmica & Superado por control digital \\
        \bottomrule
    \end{tabularx}
\end{table}

\section{Oportunidad y aportación de la propuesta}
\label{sec:arte-oportunidad}

A partir de la revisión realizada, se identifica una oportunidad de desarrollo enfocada en optimizar la relación costo-rendimiento mediante la integración de hardware de procesamiento accesible y técnicas de control robustas adaptadas a las condiciones operativas locales.
